\section{Turli muhitlarda elektr toki}

\subsection*{Metallarda elektr toki}



\begin{enumerate}

\item Metallarda elektr tokini qanday zarralar oqimi hosil qiladi?
\begin{enumerate}[label=\Alph*)]
\item elektronlar.
\item ionlar.
\item kovaklar.
\item molekulalar.
\end{enumerate}

\item Metallarda elektr toki qanday elektr zaryadlarning tartibli harakati tufayli hosil bo‘ladi?
\begin{enumerate}[label=\Alph*)]
\item musbat va manfiy ionlar, elektronlar.
\item musbat va manfiy ionlar.
\item faqat elektronlar.
\item elektron va musbat ionlar.
\end{enumerate}

\item Metallardan elektr toki o‘tganda tokning qanday ta’sirlari jarayon bo‘ladi?
\begin{enumerate}[label=\Alph*)]
\item faqat issiqlik va magnit ta’sirlari.
\item faqat issiqlik va kimyoviy ta’sirlari.
\item faqat kimyoviy va magnit ta’sirlari.
\item issiqlik, magnit va kimyoviy ta’sirlari.
\end{enumerate}

\item Mandelshtam–Papaleksi, Styuart–Tolmen tajribalari tasdiqlashicha, metallarda elektr o‘tkazuvchanlik ...
\begin{enumerate}[label=\Alph*)]
\item elektronli.
\item ionli.
\item elektron-ionli.
\item elektron-kovakli.
\end{enumerate}

\item Metallar o‘tkazuvchanligining klassik elektron nazariyasiga ko‘ra $j/en$ ifoda qaysi kattalikni beradi? ($j$ – tok zichligi, $e$ – elektron zaryadi, $n$ – elektronlarning konsentratsiyasi).
\begin{enumerate}[label=\Alph*)]
\item solishtirma qarshilik.
\item tok kuchi.
\item solishtirma o‘tkazuvchanlik.
\item elektronlar tartibli harakatining o‘rtacha tezligi.
\end{enumerate}

\item Metallarning harorati ortishi bilan elektr qarshiligi qanday o‘zgaradi?
\begin{enumerate}[label=\Alph*)]
\item chiziqli oshadi.
\item o‘zgarmaydi.
\item chiziqli kamayadi.
\item keskin oshadi.
\end{enumerate}

\item O‘tkazgich temperaturasini o‘zgarmas qoldirib, undan o‘tayotgan tok kuchini oshirsak, uning qarshiligi ...
\begin{enumerate}[label=\Alph*)]
\item oshishi ham, kamayishi ham mumkin.
\item oshadi.
\item kamayadi.
\item o‘zgarmaydi.
\end{enumerate}


\item G‘altakning mis simi 25$^\circ$C da 11 $\Omega$ qarshilikka ega. 100$^\circ$C gacha isitilganda uning qarshiligi necha $\Omega$ bo‘lishini baholang. Mis qarshiligining temperaturaviy koeffitsienti $4 \cdot 10^{-3}$ K$^{-1}$.
\begin{enumerate}[label=\Alph*)]
\item 15.
\item 14.
\item 13.
\item 12.
\end{enumerate}

\item Mis simning 50$^\circ$C dagi qarshiligi 2 $\Omega$ bo‘lsa, uning 100$^\circ$C dagi qarshiligi necha $\Omega$ bo‘ladi? Mis uchun qarshilikning temperaturaviy koeffitsienti 0,004 K$^{-1}$.
\begin{enumerate}[label=\Alph*)]
\item 0,5.
\item 4,66.
\item 1,16.
\item 2,33.
\end{enumerate}

\item 0$^\circ$C temperaturada platinali termometrning qarshiligi 6 $\Omega$, 100$^\circ$C da esa 8,4 $\Omega$ ga teng. Qanday temperaturada ($^\circ$C) uning qarshiligi 14,4 $\Omega$ ga teng bo‘ladi?
\begin{enumerate}[label=\Alph*)]
\item 300.
\item 250.
\item 200.
\item 350.
\end{enumerate}

\item Lampa tolasining 0$^\circ$C haroratdagi qarshiligi 2000$^\circ$C haroratdagi qarshiligidan 15 marta kichik. Tola qarshiligining temperaturaviy koeffitsienti qanday ($10^{-3}$ K$^{-1}$)?
\begin{enumerate}[label=\Alph*)]
\item 3,5.
\item 7.
\item 14.
\item 21.
\end{enumerate}

\item 0$^\circ$C da qarshiligi 80 $\Omega$ bo‘lgan o‘tkazgichdan tok o‘tishi natijasida uning qarshiligi 4 $\Omega$ ga o‘zgardi. Agar qarshilikning termik koeffitsienti 0,004 1/K ga teng bo‘lsa, uning temperaturasi qanchaga ($^\circ$C) o‘zgargan?
\begin{enumerate}[label=\Alph*)]
\item 12,5.
\item 50.
\item 20.
\item 25.
\end{enumerate}

\item O‘ta o‘tkazuvchanlik hodisasini kim kashf etgan?
\begin{enumerate}[label=\Alph*)]
\item X.K. Onnes.
\item N.N. Bogolyubov.
\item L.N. Kuper.
\item J. Bardin.
\end{enumerate}

\item O‘ta o‘tkazuvchanlik birinchi marta qaysi moddada kashf etilgan?
\begin{enumerate}[label=\Alph*)]
\item mis.
\item alyuminiy.
\item temir.
\item simob.
\end{enumerate}



\end{enumerate}


\begin{table}[!ht]
\centering
\begin{tabular}{|c|c|c|c|c|c|c|}
\hline
1a & 2c & 3a & 4a & 5d & 6a & 7d \\ \hline
8b & 9d & 10d & 11b & 12a & 13a & 14d \\ \hline
\end{tabular}
\end{table}


\subsection*{Yarim o`tkazgichlarda elektr toki}


\begin{enumerate}

\item Temperatura ortishi bilan yarimo‘tkazgichning qarshiligi ...
\begin{enumerate}[label=\Alph*)]
\item kamayadi.
\item o‘zgarmaydi.
\item ortadi.
\item avval ortadi, so‘ngra kamayadi.
\end{enumerate}

\item Yarimo‘tkazgichlarda elektr tokini qanday zarralar oqimi hosil qiladi? Eng to‘la javobni ko‘rsating.
\begin{enumerate}[label=\Alph*)]
\item elektronlar va kovaklar.
\item elektronlar va ionlar.
\item ionlar va kovaklar.
\item atomlar va elektronlar.
\end{enumerate}

\item Aralashmasiz yarimo‘tkazgichlar qanday turdagi o‘tkazuvchanlikka ega?
\begin{enumerate}[label=\Alph*)]
\item elektr tokini o‘tkazmaydilar.
\item teng miqdorda elektron va teshikli o‘tkazuvchanlikka.
\item asosan teshikli o‘tkazuvchanlikka.
\item asosan elektron o‘tkazuvchanlikka.
\end{enumerate}

\item Yarimo‘tkazgichda teshik va elektron uchrashganda nima hosil bo‘ladi?
\begin{enumerate}[label=\Alph*)]
\item musbat ion.
\item manfiy ion.
\item neytral atom.
\item musbat va manfiy ionlar.
\end{enumerate}

\item Yarimo‘tkazgichda elektron va teshik uchrashganda qanday hodisa yuz beradi?
\begin{enumerate}[label=\Alph*)]
\item energiya yutiladi.
\item musbat ion hosil bo‘ladi.
\item energiya ajraladi.
\item manfiy ion hosil bo‘ladi.
\end{enumerate}

\item Toza yarimo‘tkazgichdan elektronlarning tartibli harakati tufayli 1,7 mA tok o‘tmoqda. Yarimo‘tkazgichdan o‘tayotgan to‘la tok kuchi qanday (mA)?
\begin{enumerate}[label=\Alph*)]
\item 0.
\item 0,85.
\item 1,7.
\item 3,4.
\end{enumerate}

\item Yarimo‘tkazgichlarda donor aralashmasi bo‘lganda, ... ortadi. Mazkur jumladagi nuqtalar o‘rnini to‘ldiring.
\begin{enumerate}[label=\Alph*)]
\item kovaklar soni.
\item erkin pozitronlar soni.
\item erkin protonlar soni.
\item erkin elektronlar soni.
\end{enumerate}

\item Yarimo‘tkazgichlarda akseptor aralashmasi bo‘lganda, ... ortadi. Mazkur jumladagi nuqtalar o‘rnini to‘ldiring.
\begin{enumerate}[label=\Alph*)]
\item erkin protonlar soni.
\item erkin elektronlar soni.
\item kovaklar soni.
\item erkin pozitronlar soni.
\end{enumerate}

\item Yarimo‘tkazgich asosan teshikli o‘tkazuvchanlikka ega. Kristallda qanday aralashma bor?
\begin{enumerate}[label=\Alph*)]
\item asosan teshikli o‘tkazuvchanlikka ega yarimo‘tkazgich bo‘lmaydi.
\item akseptor aralashma.
\item aralashmalar yo‘q.
\item donor aralashma.
\end{enumerate}

\item Yarimo‘tkazgich asosan elektron o‘tkazuvchanlikka ega. Kristallda qanday aralashma bor?
\begin{enumerate}[label=\Alph*)]
\item donor aralashma.
\item akseptor aralashma.
\item aralashmalar yo‘q.
\item donor va akseptorlarning konsentratsiyalari teng.
\end{enumerate}

\item Donor aralashmali yarimo‘tkazgichlar qaysi turdagi o‘tkazuvchanlikka ega?
\begin{enumerate}[label=\Alph*)]
\item asosan elektronli.
\item asosan kovakli.
\item ham kovakli, ham elektronli (teng miqdorda).
\item ionli.
\end{enumerate}

\item Akseptor aralashmali yarimo‘tkazgichlar qanday tipdagi o‘tkazuvchanlikka ega?
\begin{enumerate}[label=\Alph*)]
\item asosan teshikli o‘tkazuvchanlikka.
\item asosan elektron o‘tkazuvchanlikka.
\item teng miqdorda elektron va teshikli o‘tkazuvchanlikka.
\item elektr tokini o‘tkazmaydi.
\end{enumerate}

\item Quyidagi lampalarning qaysi biri yonadi? (Rasmga qarang)

\begin{circuitikz}[european, scale=0.8]
% 1-chizma: O'zgarmas tok va kondensator
\draw (0,0) to[battery1] (0,2) to[lamp, l=1)] (2,2) to[C] (2,0) -- (0,0);

% 2-chizma: O'zgaruvchan tok va kondensator
\begin{scope}[xshift=4cm]
\draw (0,0) to[vsourcesin] (0,2) to[lamp, l=2)] (2,2) to[C] (2,0) -- (0,0);
\end{scope}

% 3-chizma: Diod (teskari ulanish)
\begin{scope}[yshift=-3.5cm]
\draw (0,0) to[battery1] (0,2) to[lamp, l=3)] (2,2) to[full diode, invert] (2,0) -- (0,0);
\end{scope}

% 4-chizma: Diod (to'g'ri ulanish)
\begin{scope}[xshift=4cm, yshift=-3.5cm]
\draw (0,0) to[battery1] (0,2) to[lamp, l=4)] (2,2) to[full diode] (2,0) -- (0,0);
\end{scope}
\end{circuitikz}

\begin{enumerate}[label=\Alph*)]
\item 2, 3.
\item 2, 4.
\item 2.
\item 3.
\end{enumerate}

\item Tranzistor qanday vazifani bajaradi?
\begin{enumerate}[label=\Alph*)]
\item o‘zgarmas tokni o‘zgaruvchan tokka aylantiradi.
\item o‘zgaruvchan tokni o‘zgarmas tokka aylantiradi.
\item elektr signallarini kuchaytiradi.
\item elektr zaryadlarini to‘playdi.
\end{enumerate}

\item Tranzistor nimalardan tashkil topgan?
\begin{enumerate}[label=\Alph*)]
\item emitter, baza va kollektordan.
\item emitter va kollektordan.
\item baza va emitterdan.
\item anod, katod va to‘rdan.
\end{enumerate}

\end{enumerate}


\begin{table}[!ht]
\centering
\begin{tabular}{|c|c|c|c|c|c|c|c|c|c|}
\hline
1a & 2a & 3b & 4c & 5c & 6d & 7d & 8c & 9b & 10a \\ \hline
11a & 12a & 13a & 14c & 15a & & & & & \\ \hline
\end{tabular}
\end{table}





\subsection*{Vakumda elektr toki}


\begin{enumerate}

\item Termoelektron emissiya nima?
\begin{enumerate}[label=\Alph*)]
\item qizdirilgan qattiq jism sirtidan musbat ionlarning nurlanishi.
\item atom yadrosining parchalanishi.
\item atomning elektron va ionga ajralishi.
\item qizdirilgan qattiq jism sirtidan elektronlarning nurlanishi.
\end{enumerate}

\item Katoddan birlik vaqtda bug‘lanib chiqayotgan elektronlar soni nimalarga bog‘liq? Eng to‘liq javobni belgilang.
\begin{enumerate}[label=\Alph*)]
\item katod materialiga, haroratga.
\item katod materialiga, kuchlanishga.
\item haroratga, sirt o‘lchamlariga, tok kuchiga.
\item katod materialiga, sirt o‘lchamlariga, haroratga.
\end{enumerate}

\item Agar katoddan har sekundda $5 \cdot 10^{16}$ ta elektron ajralib chiqayotgan bo‘lsa, anod zanjiridagi tok qanday (mA)?
\begin{enumerate}[label=\Alph*)]
\item 5.
\item 8.
\item 9.
\item 10.
\end{enumerate}

\item Vakuumda elektr toki qanday zarrachalarning harakati tufayli vujudga keladi?
\begin{enumerate}[label=\Alph*)]
\item protonlar.
\item ixtiyoriy zaryadlangan zarrachalar.
\item ionlar.
\item elektronlar.
\end{enumerate}

\item Katod nurlari – katoddan anodga tez harakat qiluvchi ... oqimi. Mazkur jumladagi nuqtalar o‘rnini to‘ldiring.
\begin{enumerate}[label=\Alph*)]
\item pozitronlar.
\item protonlar.
\item ionlar.
\item elektronlar.
\end{enumerate}


\item Vakuumli dioddagi tok kuchi to‘yinishga erishmagan. Agar diodning anod kuchlanishi 4 marta ortsa, tok kuchi qanday o‘zgaradi?
\begin{enumerate}[label=\Alph*)]
\item o‘zgarmaydi.
\item 4 marta ortadi.
\item 8 marta ortadi.
\item 16 marta ortadi.
\end{enumerate}

\item Vakuumli dioddagi tok kuchi to‘yinishga erishmagan. Diodda ajralayotgan quvvat qaysi formulada to‘g‘ri ifodalangan?
\begin{enumerate}[label=\Alph*)]
\item $P = kU^{\frac{3}{2}}$.
\item $P = kU^{\frac{1}{2}}$.
\item $P = kU^{\frac{5}{2}}$.
\item $P = kU^2$.
\end{enumerate}

\item Vakuumli dioddagi tok kuchi to‘yinishga erishmagan. Agar diodning anod kuchlanishi 4 marta ortsa, ajralayotgan quvvat qanday o‘zgaradi?
\begin{enumerate}[label=\Alph*)]
\item o‘zgarmaydi.
\item 4 marta ortadi.
\item 8 marta ortadi.
\item 32 marta ortadi.
\end{enumerate}

\item Qaysi asboblar o‘zgaruvchan tokni o‘zgarmas tokka aylantirishda asosiy element bo‘ladi? 1) transformator; 2) yarimo‘tkazgichli diod; 3) elektron lampa–diod; 4) tranzistor.
\begin{enumerate}[label=\Alph*)]
\item 2, 3.
\item 2, 4.
\item 1, 2.
\item 1, 2, 3.
\end{enumerate}

\item Vakuumli diodda katoddan chiqib, anodga yetib borgan elektronlarning kinetik energiyasini hisoblash uchun nimalarni bilish kerak? 1) elektron zaryadini; 2) anod va katod orasidagi potensiallar ayirmasini; 3) anod va katod orasidagi masofani.
\begin{enumerate}[label=\Alph*)]
\item 1.
\item 2.
\item 1, 2, 3.
\item 1, 2.
\end{enumerate}

\item Anod bilan katod orasidagi potensiallar farqi 360 V. Ular orasidagi masofa 1 cm bo‘lsa, elektron necha m/s$^2$ tezlanish oladi? $m=9,1 \cdot 10^{-31}$ kg, $e=1,6 \cdot 10^{-19}$ C.
\begin{enumerate}[label=\Alph*)]
\item $6,3 \cdot 10^{15}$.
\item $6 \cdot 10^{14}$.
\item $5 \cdot 10^{15}$.
\item $1,6 \cdot 10^{15}$.
\end{enumerate}


\item Elektr lampadagi tok kuchi 1,28 A ga teng. Uning spiralining ko‘ndalang kesimi orqali 5 minutda nechta elektron o‘tadi? $e=1,6 \cdot 10^{-19}$ C.
\begin{enumerate}[label=\Alph*)]
\item $2,5 \cdot 10^{20}$.
\item $2,4 \cdot 10^{20}$.
\item $2,4 \cdot 10^{22}$.
\item $2,4 \cdot 10^{21}$.
\end{enumerate}

\item Diodda elektron anodga $8 \cdot 10^6$ m/s tezlik bilan yetib keldi. Agar elektronning katod yaqinidagi tezligi nol bo‘lsa, anod kuchlanishini toping (V).
\begin{enumerate}[label=\Alph*)]
\item 220.
\item 160.
\item 182.
\item 280.
\end{enumerate}

\item Elektron-nur trubkasining anodi va katodi orasidagi potensiallar ayirmasi $U$ ga teng bo‘lsa, elektron erishadigan tezlik qanday? $m$ va $e$ – elektronning massasi va zaryadi.
\begin{enumerate}[label=\Alph*)]
\item $\sqrt{2emU}$.
\item $\sqrt{\frac{eU}{m}}$.
\item $\sqrt{\frac{2eU}{m}}$.
\item $\sqrt{\frac{2mU}{e}}$.
\end{enumerate}

\item Televizor trubkasida anoddan ekrangacha bo‘lgan masofani elektronlar 4 ns da o‘tadi. Kineskopdagi tezlashtiruvchi kuchlanish necha kV? Anoddan ekrangacha bo‘lgan masofa 30 cm. Boshlang‘ich tezlik nol.
\begin{enumerate}[label=\Alph*)]
\item 1,6.
\item 2.
\item 3.
\item 16.
\end{enumerate}

\end{enumerate}


\begin{table}[!ht]
\centering
\begin{tabular}{|c|c|c|c|c|c|c|c|c|c|}
\hline
1d & 2d & 3b & 4d & 5d & 6c & 7c & 8d & 9a & 10d \\ \hline
11a & 12d & 13c & 14c & 15d & & & & & \\ \hline
\end{tabular}
\end{table}

